% gerar-fundos.tex — gera fundo-impar.pdf, fundo-par.pdf e capa.pdf (A5) dos
% MWEs web e print. Arte provisória: faixas de cor, sem texto.
% Uso, uma vez por arquivo (o nome do PDF é o do -jobname):
%   lualatex -jobname=fundo-impar '\def\qual{impar}% gerar-fundos.tex — gera fundo-impar.pdf, fundo-par.pdf e capa.pdf (A5) dos
% MWEs web e print. Arte provisória: faixas de cor, sem texto.
% Uso, uma vez por arquivo (o nome do PDF é o do -jobname):
%   lualatex -jobname=fundo-impar '\def\qual{impar}% gerar-fundos.tex — gera fundo-impar.pdf, fundo-par.pdf e capa.pdf (A5) dos
% MWEs web e print. Arte provisória: faixas de cor, sem texto.
% Uso, uma vez por arquivo (o nome do PDF é o do -jobname):
%   lualatex -jobname=fundo-impar '\def\qual{impar}% gerar-fundos.tex — gera fundo-impar.pdf, fundo-par.pdf e capa.pdf (A5) dos
% MWEs web e print. Arte provisória: faixas de cor, sem texto.
% Uso, uma vez por arquivo (o nome do PDF é o do -jobname):
%   lualatex -jobname=fundo-impar '\def\qual{impar}\input{gerar-fundos}'
%   lualatex -jobname=fundo-par   '\def\qual{par}\input{gerar-fundos}'
%   lualatex -jobname=capa        '\def\qual{capa}\input{gerar-fundos}'
% Os PDFs gerados ficam versionados ao lado (pequenos, vetoriais).
\documentclass{article}
\usepackage[paperwidth=148mm,paperheight=210mm,margin=0pt]{geometry}
\usepackage{xcolor,eso-pic}
\pagestyle{empty}
\def\faixa#1#2#3#4#5{\put(#1,#2){\color[RGB]{#5}\rule{#3}{#4}}}
\begin{document}
\AddToShipoutPictureBG*{\AtPageLowerLeft{%
  \expandafter\ifx\csname qual\endcsname\relax\else
  \def\a{impar}\ifx\qual\a   % faixa na margem externa à direita
    \faixa{130mm}{0mm}{18mm}{210mm}{217,238,242}\faixa{0mm}{0mm}{148mm}{3mm}{4,88,110}\fi
  \def\a{par}\ifx\qual\a     % faixa na margem externa à esquerda
    \faixa{0mm}{0mm}{18mm}{210mm}{217,238,242}\faixa{0mm}{0mm}{148mm}{3mm}{4,88,110}\fi
  \def\a{capa}\ifx\qual\a
    \faixa{0mm}{0mm}{148mm}{210mm}{25,159,175}\faixa{0mm}{150mm}{148mm}{20mm}{191,115,171}\fi
  \fi}}
\null
\end{document}
'
%   lualatex -jobname=fundo-par   '\def\qual{par}% gerar-fundos.tex — gera fundo-impar.pdf, fundo-par.pdf e capa.pdf (A5) dos
% MWEs web e print. Arte provisória: faixas de cor, sem texto.
% Uso, uma vez por arquivo (o nome do PDF é o do -jobname):
%   lualatex -jobname=fundo-impar '\def\qual{impar}\input{gerar-fundos}'
%   lualatex -jobname=fundo-par   '\def\qual{par}\input{gerar-fundos}'
%   lualatex -jobname=capa        '\def\qual{capa}\input{gerar-fundos}'
% Os PDFs gerados ficam versionados ao lado (pequenos, vetoriais).
\documentclass{article}
\usepackage[paperwidth=148mm,paperheight=210mm,margin=0pt]{geometry}
\usepackage{xcolor,eso-pic}
\pagestyle{empty}
\def\faixa#1#2#3#4#5{\put(#1,#2){\color[RGB]{#5}\rule{#3}{#4}}}
\begin{document}
\AddToShipoutPictureBG*{\AtPageLowerLeft{%
  \expandafter\ifx\csname qual\endcsname\relax\else
  \def\a{impar}\ifx\qual\a   % faixa na margem externa à direita
    \faixa{130mm}{0mm}{18mm}{210mm}{217,238,242}\faixa{0mm}{0mm}{148mm}{3mm}{4,88,110}\fi
  \def\a{par}\ifx\qual\a     % faixa na margem externa à esquerda
    \faixa{0mm}{0mm}{18mm}{210mm}{217,238,242}\faixa{0mm}{0mm}{148mm}{3mm}{4,88,110}\fi
  \def\a{capa}\ifx\qual\a
    \faixa{0mm}{0mm}{148mm}{210mm}{25,159,175}\faixa{0mm}{150mm}{148mm}{20mm}{191,115,171}\fi
  \fi}}
\null
\end{document}
'
%   lualatex -jobname=capa        '\def\qual{capa}% gerar-fundos.tex — gera fundo-impar.pdf, fundo-par.pdf e capa.pdf (A5) dos
% MWEs web e print. Arte provisória: faixas de cor, sem texto.
% Uso, uma vez por arquivo (o nome do PDF é o do -jobname):
%   lualatex -jobname=fundo-impar '\def\qual{impar}\input{gerar-fundos}'
%   lualatex -jobname=fundo-par   '\def\qual{par}\input{gerar-fundos}'
%   lualatex -jobname=capa        '\def\qual{capa}\input{gerar-fundos}'
% Os PDFs gerados ficam versionados ao lado (pequenos, vetoriais).
\documentclass{article}
\usepackage[paperwidth=148mm,paperheight=210mm,margin=0pt]{geometry}
\usepackage{xcolor,eso-pic}
\pagestyle{empty}
\def\faixa#1#2#3#4#5{\put(#1,#2){\color[RGB]{#5}\rule{#3}{#4}}}
\begin{document}
\AddToShipoutPictureBG*{\AtPageLowerLeft{%
  \expandafter\ifx\csname qual\endcsname\relax\else
  \def\a{impar}\ifx\qual\a   % faixa na margem externa à direita
    \faixa{130mm}{0mm}{18mm}{210mm}{217,238,242}\faixa{0mm}{0mm}{148mm}{3mm}{4,88,110}\fi
  \def\a{par}\ifx\qual\a     % faixa na margem externa à esquerda
    \faixa{0mm}{0mm}{18mm}{210mm}{217,238,242}\faixa{0mm}{0mm}{148mm}{3mm}{4,88,110}\fi
  \def\a{capa}\ifx\qual\a
    \faixa{0mm}{0mm}{148mm}{210mm}{25,159,175}\faixa{0mm}{150mm}{148mm}{20mm}{191,115,171}\fi
  \fi}}
\null
\end{document}
'
% Os PDFs gerados ficam versionados ao lado (pequenos, vetoriais).
\documentclass{article}
\usepackage[paperwidth=148mm,paperheight=210mm,margin=0pt]{geometry}
\usepackage{xcolor,eso-pic}
\pagestyle{empty}
\def\faixa#1#2#3#4#5{\put(#1,#2){\color[RGB]{#5}\rule{#3}{#4}}}
\begin{document}
\AddToShipoutPictureBG*{\AtPageLowerLeft{%
  \expandafter\ifx\csname qual\endcsname\relax\else
  \def\a{impar}\ifx\qual\a   % faixa na margem externa à direita
    \faixa{130mm}{0mm}{18mm}{210mm}{217,238,242}\faixa{0mm}{0mm}{148mm}{3mm}{4,88,110}\fi
  \def\a{par}\ifx\qual\a     % faixa na margem externa à esquerda
    \faixa{0mm}{0mm}{18mm}{210mm}{217,238,242}\faixa{0mm}{0mm}{148mm}{3mm}{4,88,110}\fi
  \def\a{capa}\ifx\qual\a
    \faixa{0mm}{0mm}{148mm}{210mm}{25,159,175}\faixa{0mm}{150mm}{148mm}{20mm}{191,115,171}\fi
  \fi}}
\null
\end{document}
'
%   lualatex -jobname=fundo-par   '\def\qual{par}% gerar-fundos.tex — gera fundo-impar.pdf, fundo-par.pdf e capa.pdf (A5) dos
% MWEs web e print. Arte provisória: faixas de cor, sem texto.
% Uso, uma vez por arquivo (o nome do PDF é o do -jobname):
%   lualatex -jobname=fundo-impar '\def\qual{impar}% gerar-fundos.tex — gera fundo-impar.pdf, fundo-par.pdf e capa.pdf (A5) dos
% MWEs web e print. Arte provisória: faixas de cor, sem texto.
% Uso, uma vez por arquivo (o nome do PDF é o do -jobname):
%   lualatex -jobname=fundo-impar '\def\qual{impar}\input{gerar-fundos}'
%   lualatex -jobname=fundo-par   '\def\qual{par}\input{gerar-fundos}'
%   lualatex -jobname=capa        '\def\qual{capa}\input{gerar-fundos}'
% Os PDFs gerados ficam versionados ao lado (pequenos, vetoriais).
\documentclass{article}
\usepackage[paperwidth=148mm,paperheight=210mm,margin=0pt]{geometry}
\usepackage{xcolor,eso-pic}
\pagestyle{empty}
\def\faixa#1#2#3#4#5{\put(#1,#2){\color[RGB]{#5}\rule{#3}{#4}}}
\begin{document}
\AddToShipoutPictureBG*{\AtPageLowerLeft{%
  \expandafter\ifx\csname qual\endcsname\relax\else
  \def\a{impar}\ifx\qual\a   % faixa na margem externa à direita
    \faixa{130mm}{0mm}{18mm}{210mm}{217,238,242}\faixa{0mm}{0mm}{148mm}{3mm}{4,88,110}\fi
  \def\a{par}\ifx\qual\a     % faixa na margem externa à esquerda
    \faixa{0mm}{0mm}{18mm}{210mm}{217,238,242}\faixa{0mm}{0mm}{148mm}{3mm}{4,88,110}\fi
  \def\a{capa}\ifx\qual\a
    \faixa{0mm}{0mm}{148mm}{210mm}{25,159,175}\faixa{0mm}{150mm}{148mm}{20mm}{191,115,171}\fi
  \fi}}
\null
\end{document}
'
%   lualatex -jobname=fundo-par   '\def\qual{par}% gerar-fundos.tex — gera fundo-impar.pdf, fundo-par.pdf e capa.pdf (A5) dos
% MWEs web e print. Arte provisória: faixas de cor, sem texto.
% Uso, uma vez por arquivo (o nome do PDF é o do -jobname):
%   lualatex -jobname=fundo-impar '\def\qual{impar}\input{gerar-fundos}'
%   lualatex -jobname=fundo-par   '\def\qual{par}\input{gerar-fundos}'
%   lualatex -jobname=capa        '\def\qual{capa}\input{gerar-fundos}'
% Os PDFs gerados ficam versionados ao lado (pequenos, vetoriais).
\documentclass{article}
\usepackage[paperwidth=148mm,paperheight=210mm,margin=0pt]{geometry}
\usepackage{xcolor,eso-pic}
\pagestyle{empty}
\def\faixa#1#2#3#4#5{\put(#1,#2){\color[RGB]{#5}\rule{#3}{#4}}}
\begin{document}
\AddToShipoutPictureBG*{\AtPageLowerLeft{%
  \expandafter\ifx\csname qual\endcsname\relax\else
  \def\a{impar}\ifx\qual\a   % faixa na margem externa à direita
    \faixa{130mm}{0mm}{18mm}{210mm}{217,238,242}\faixa{0mm}{0mm}{148mm}{3mm}{4,88,110}\fi
  \def\a{par}\ifx\qual\a     % faixa na margem externa à esquerda
    \faixa{0mm}{0mm}{18mm}{210mm}{217,238,242}\faixa{0mm}{0mm}{148mm}{3mm}{4,88,110}\fi
  \def\a{capa}\ifx\qual\a
    \faixa{0mm}{0mm}{148mm}{210mm}{25,159,175}\faixa{0mm}{150mm}{148mm}{20mm}{191,115,171}\fi
  \fi}}
\null
\end{document}
'
%   lualatex -jobname=capa        '\def\qual{capa}% gerar-fundos.tex — gera fundo-impar.pdf, fundo-par.pdf e capa.pdf (A5) dos
% MWEs web e print. Arte provisória: faixas de cor, sem texto.
% Uso, uma vez por arquivo (o nome do PDF é o do -jobname):
%   lualatex -jobname=fundo-impar '\def\qual{impar}\input{gerar-fundos}'
%   lualatex -jobname=fundo-par   '\def\qual{par}\input{gerar-fundos}'
%   lualatex -jobname=capa        '\def\qual{capa}\input{gerar-fundos}'
% Os PDFs gerados ficam versionados ao lado (pequenos, vetoriais).
\documentclass{article}
\usepackage[paperwidth=148mm,paperheight=210mm,margin=0pt]{geometry}
\usepackage{xcolor,eso-pic}
\pagestyle{empty}
\def\faixa#1#2#3#4#5{\put(#1,#2){\color[RGB]{#5}\rule{#3}{#4}}}
\begin{document}
\AddToShipoutPictureBG*{\AtPageLowerLeft{%
  \expandafter\ifx\csname qual\endcsname\relax\else
  \def\a{impar}\ifx\qual\a   % faixa na margem externa à direita
    \faixa{130mm}{0mm}{18mm}{210mm}{217,238,242}\faixa{0mm}{0mm}{148mm}{3mm}{4,88,110}\fi
  \def\a{par}\ifx\qual\a     % faixa na margem externa à esquerda
    \faixa{0mm}{0mm}{18mm}{210mm}{217,238,242}\faixa{0mm}{0mm}{148mm}{3mm}{4,88,110}\fi
  \def\a{capa}\ifx\qual\a
    \faixa{0mm}{0mm}{148mm}{210mm}{25,159,175}\faixa{0mm}{150mm}{148mm}{20mm}{191,115,171}\fi
  \fi}}
\null
\end{document}
'
% Os PDFs gerados ficam versionados ao lado (pequenos, vetoriais).
\documentclass{article}
\usepackage[paperwidth=148mm,paperheight=210mm,margin=0pt]{geometry}
\usepackage{xcolor,eso-pic}
\pagestyle{empty}
\def\faixa#1#2#3#4#5{\put(#1,#2){\color[RGB]{#5}\rule{#3}{#4}}}
\begin{document}
\AddToShipoutPictureBG*{\AtPageLowerLeft{%
  \expandafter\ifx\csname qual\endcsname\relax\else
  \def\a{impar}\ifx\qual\a   % faixa na margem externa à direita
    \faixa{130mm}{0mm}{18mm}{210mm}{217,238,242}\faixa{0mm}{0mm}{148mm}{3mm}{4,88,110}\fi
  \def\a{par}\ifx\qual\a     % faixa na margem externa à esquerda
    \faixa{0mm}{0mm}{18mm}{210mm}{217,238,242}\faixa{0mm}{0mm}{148mm}{3mm}{4,88,110}\fi
  \def\a{capa}\ifx\qual\a
    \faixa{0mm}{0mm}{148mm}{210mm}{25,159,175}\faixa{0mm}{150mm}{148mm}{20mm}{191,115,171}\fi
  \fi}}
\null
\end{document}
'
%   lualatex -jobname=capa        '\def\qual{capa}% gerar-fundos.tex — gera fundo-impar.pdf, fundo-par.pdf e capa.pdf (A5) dos
% MWEs web e print. Arte provisória: faixas de cor, sem texto.
% Uso, uma vez por arquivo (o nome do PDF é o do -jobname):
%   lualatex -jobname=fundo-impar '\def\qual{impar}% gerar-fundos.tex — gera fundo-impar.pdf, fundo-par.pdf e capa.pdf (A5) dos
% MWEs web e print. Arte provisória: faixas de cor, sem texto.
% Uso, uma vez por arquivo (o nome do PDF é o do -jobname):
%   lualatex -jobname=fundo-impar '\def\qual{impar}\input{gerar-fundos}'
%   lualatex -jobname=fundo-par   '\def\qual{par}\input{gerar-fundos}'
%   lualatex -jobname=capa        '\def\qual{capa}\input{gerar-fundos}'
% Os PDFs gerados ficam versionados ao lado (pequenos, vetoriais).
\documentclass{article}
\usepackage[paperwidth=148mm,paperheight=210mm,margin=0pt]{geometry}
\usepackage{xcolor,eso-pic}
\pagestyle{empty}
\def\faixa#1#2#3#4#5{\put(#1,#2){\color[RGB]{#5}\rule{#3}{#4}}}
\begin{document}
\AddToShipoutPictureBG*{\AtPageLowerLeft{%
  \expandafter\ifx\csname qual\endcsname\relax\else
  \def\a{impar}\ifx\qual\a   % faixa na margem externa à direita
    \faixa{130mm}{0mm}{18mm}{210mm}{217,238,242}\faixa{0mm}{0mm}{148mm}{3mm}{4,88,110}\fi
  \def\a{par}\ifx\qual\a     % faixa na margem externa à esquerda
    \faixa{0mm}{0mm}{18mm}{210mm}{217,238,242}\faixa{0mm}{0mm}{148mm}{3mm}{4,88,110}\fi
  \def\a{capa}\ifx\qual\a
    \faixa{0mm}{0mm}{148mm}{210mm}{25,159,175}\faixa{0mm}{150mm}{148mm}{20mm}{191,115,171}\fi
  \fi}}
\null
\end{document}
'
%   lualatex -jobname=fundo-par   '\def\qual{par}% gerar-fundos.tex — gera fundo-impar.pdf, fundo-par.pdf e capa.pdf (A5) dos
% MWEs web e print. Arte provisória: faixas de cor, sem texto.
% Uso, uma vez por arquivo (o nome do PDF é o do -jobname):
%   lualatex -jobname=fundo-impar '\def\qual{impar}\input{gerar-fundos}'
%   lualatex -jobname=fundo-par   '\def\qual{par}\input{gerar-fundos}'
%   lualatex -jobname=capa        '\def\qual{capa}\input{gerar-fundos}'
% Os PDFs gerados ficam versionados ao lado (pequenos, vetoriais).
\documentclass{article}
\usepackage[paperwidth=148mm,paperheight=210mm,margin=0pt]{geometry}
\usepackage{xcolor,eso-pic}
\pagestyle{empty}
\def\faixa#1#2#3#4#5{\put(#1,#2){\color[RGB]{#5}\rule{#3}{#4}}}
\begin{document}
\AddToShipoutPictureBG*{\AtPageLowerLeft{%
  \expandafter\ifx\csname qual\endcsname\relax\else
  \def\a{impar}\ifx\qual\a   % faixa na margem externa à direita
    \faixa{130mm}{0mm}{18mm}{210mm}{217,238,242}\faixa{0mm}{0mm}{148mm}{3mm}{4,88,110}\fi
  \def\a{par}\ifx\qual\a     % faixa na margem externa à esquerda
    \faixa{0mm}{0mm}{18mm}{210mm}{217,238,242}\faixa{0mm}{0mm}{148mm}{3mm}{4,88,110}\fi
  \def\a{capa}\ifx\qual\a
    \faixa{0mm}{0mm}{148mm}{210mm}{25,159,175}\faixa{0mm}{150mm}{148mm}{20mm}{191,115,171}\fi
  \fi}}
\null
\end{document}
'
%   lualatex -jobname=capa        '\def\qual{capa}% gerar-fundos.tex — gera fundo-impar.pdf, fundo-par.pdf e capa.pdf (A5) dos
% MWEs web e print. Arte provisória: faixas de cor, sem texto.
% Uso, uma vez por arquivo (o nome do PDF é o do -jobname):
%   lualatex -jobname=fundo-impar '\def\qual{impar}\input{gerar-fundos}'
%   lualatex -jobname=fundo-par   '\def\qual{par}\input{gerar-fundos}'
%   lualatex -jobname=capa        '\def\qual{capa}\input{gerar-fundos}'
% Os PDFs gerados ficam versionados ao lado (pequenos, vetoriais).
\documentclass{article}
\usepackage[paperwidth=148mm,paperheight=210mm,margin=0pt]{geometry}
\usepackage{xcolor,eso-pic}
\pagestyle{empty}
\def\faixa#1#2#3#4#5{\put(#1,#2){\color[RGB]{#5}\rule{#3}{#4}}}
\begin{document}
\AddToShipoutPictureBG*{\AtPageLowerLeft{%
  \expandafter\ifx\csname qual\endcsname\relax\else
  \def\a{impar}\ifx\qual\a   % faixa na margem externa à direita
    \faixa{130mm}{0mm}{18mm}{210mm}{217,238,242}\faixa{0mm}{0mm}{148mm}{3mm}{4,88,110}\fi
  \def\a{par}\ifx\qual\a     % faixa na margem externa à esquerda
    \faixa{0mm}{0mm}{18mm}{210mm}{217,238,242}\faixa{0mm}{0mm}{148mm}{3mm}{4,88,110}\fi
  \def\a{capa}\ifx\qual\a
    \faixa{0mm}{0mm}{148mm}{210mm}{25,159,175}\faixa{0mm}{150mm}{148mm}{20mm}{191,115,171}\fi
  \fi}}
\null
\end{document}
'
% Os PDFs gerados ficam versionados ao lado (pequenos, vetoriais).
\documentclass{article}
\usepackage[paperwidth=148mm,paperheight=210mm,margin=0pt]{geometry}
\usepackage{xcolor,eso-pic}
\pagestyle{empty}
\def\faixa#1#2#3#4#5{\put(#1,#2){\color[RGB]{#5}\rule{#3}{#4}}}
\begin{document}
\AddToShipoutPictureBG*{\AtPageLowerLeft{%
  \expandafter\ifx\csname qual\endcsname\relax\else
  \def\a{impar}\ifx\qual\a   % faixa na margem externa à direita
    \faixa{130mm}{0mm}{18mm}{210mm}{217,238,242}\faixa{0mm}{0mm}{148mm}{3mm}{4,88,110}\fi
  \def\a{par}\ifx\qual\a     % faixa na margem externa à esquerda
    \faixa{0mm}{0mm}{18mm}{210mm}{217,238,242}\faixa{0mm}{0mm}{148mm}{3mm}{4,88,110}\fi
  \def\a{capa}\ifx\qual\a
    \faixa{0mm}{0mm}{148mm}{210mm}{25,159,175}\faixa{0mm}{150mm}{148mm}{20mm}{191,115,171}\fi
  \fi}}
\null
\end{document}
'
% Os PDFs gerados ficam versionados ao lado (pequenos, vetoriais).
\documentclass{article}
\usepackage[paperwidth=148mm,paperheight=210mm,margin=0pt]{geometry}
\usepackage{xcolor,eso-pic}
\pagestyle{empty}
\def\faixa#1#2#3#4#5{\put(#1,#2){\color[RGB]{#5}\rule{#3}{#4}}}
\begin{document}
\AddToShipoutPictureBG*{\AtPageLowerLeft{%
  \expandafter\ifx\csname qual\endcsname\relax\else
  \def\a{impar}\ifx\qual\a   % faixa na margem externa à direita
    \faixa{130mm}{0mm}{18mm}{210mm}{217,238,242}\faixa{0mm}{0mm}{148mm}{3mm}{4,88,110}\fi
  \def\a{par}\ifx\qual\a     % faixa na margem externa à esquerda
    \faixa{0mm}{0mm}{18mm}{210mm}{217,238,242}\faixa{0mm}{0mm}{148mm}{3mm}{4,88,110}\fi
  \def\a{capa}\ifx\qual\a
    \faixa{0mm}{0mm}{148mm}{210mm}{25,159,175}\faixa{0mm}{150mm}{148mm}{20mm}{191,115,171}\fi
  \fi}}
\null
\end{document}
'
%   lualatex -jobname=fundo-par   '\def\qual{par}% gerar-fundos.tex — gera fundo-impar.pdf, fundo-par.pdf e capa.pdf (A5) dos
% MWEs web e print. Arte provisória: faixas de cor, sem texto.
% Uso, uma vez por arquivo (o nome do PDF é o do -jobname):
%   lualatex -jobname=fundo-impar '\def\qual{impar}% gerar-fundos.tex — gera fundo-impar.pdf, fundo-par.pdf e capa.pdf (A5) dos
% MWEs web e print. Arte provisória: faixas de cor, sem texto.
% Uso, uma vez por arquivo (o nome do PDF é o do -jobname):
%   lualatex -jobname=fundo-impar '\def\qual{impar}% gerar-fundos.tex — gera fundo-impar.pdf, fundo-par.pdf e capa.pdf (A5) dos
% MWEs web e print. Arte provisória: faixas de cor, sem texto.
% Uso, uma vez por arquivo (o nome do PDF é o do -jobname):
%   lualatex -jobname=fundo-impar '\def\qual{impar}\input{gerar-fundos}'
%   lualatex -jobname=fundo-par   '\def\qual{par}\input{gerar-fundos}'
%   lualatex -jobname=capa        '\def\qual{capa}\input{gerar-fundos}'
% Os PDFs gerados ficam versionados ao lado (pequenos, vetoriais).
\documentclass{article}
\usepackage[paperwidth=148mm,paperheight=210mm,margin=0pt]{geometry}
\usepackage{xcolor,eso-pic}
\pagestyle{empty}
\def\faixa#1#2#3#4#5{\put(#1,#2){\color[RGB]{#5}\rule{#3}{#4}}}
\begin{document}
\AddToShipoutPictureBG*{\AtPageLowerLeft{%
  \expandafter\ifx\csname qual\endcsname\relax\else
  \def\a{impar}\ifx\qual\a   % faixa na margem externa à direita
    \faixa{130mm}{0mm}{18mm}{210mm}{217,238,242}\faixa{0mm}{0mm}{148mm}{3mm}{4,88,110}\fi
  \def\a{par}\ifx\qual\a     % faixa na margem externa à esquerda
    \faixa{0mm}{0mm}{18mm}{210mm}{217,238,242}\faixa{0mm}{0mm}{148mm}{3mm}{4,88,110}\fi
  \def\a{capa}\ifx\qual\a
    \faixa{0mm}{0mm}{148mm}{210mm}{25,159,175}\faixa{0mm}{150mm}{148mm}{20mm}{191,115,171}\fi
  \fi}}
\null
\end{document}
'
%   lualatex -jobname=fundo-par   '\def\qual{par}% gerar-fundos.tex — gera fundo-impar.pdf, fundo-par.pdf e capa.pdf (A5) dos
% MWEs web e print. Arte provisória: faixas de cor, sem texto.
% Uso, uma vez por arquivo (o nome do PDF é o do -jobname):
%   lualatex -jobname=fundo-impar '\def\qual{impar}\input{gerar-fundos}'
%   lualatex -jobname=fundo-par   '\def\qual{par}\input{gerar-fundos}'
%   lualatex -jobname=capa        '\def\qual{capa}\input{gerar-fundos}'
% Os PDFs gerados ficam versionados ao lado (pequenos, vetoriais).
\documentclass{article}
\usepackage[paperwidth=148mm,paperheight=210mm,margin=0pt]{geometry}
\usepackage{xcolor,eso-pic}
\pagestyle{empty}
\def\faixa#1#2#3#4#5{\put(#1,#2){\color[RGB]{#5}\rule{#3}{#4}}}
\begin{document}
\AddToShipoutPictureBG*{\AtPageLowerLeft{%
  \expandafter\ifx\csname qual\endcsname\relax\else
  \def\a{impar}\ifx\qual\a   % faixa na margem externa à direita
    \faixa{130mm}{0mm}{18mm}{210mm}{217,238,242}\faixa{0mm}{0mm}{148mm}{3mm}{4,88,110}\fi
  \def\a{par}\ifx\qual\a     % faixa na margem externa à esquerda
    \faixa{0mm}{0mm}{18mm}{210mm}{217,238,242}\faixa{0mm}{0mm}{148mm}{3mm}{4,88,110}\fi
  \def\a{capa}\ifx\qual\a
    \faixa{0mm}{0mm}{148mm}{210mm}{25,159,175}\faixa{0mm}{150mm}{148mm}{20mm}{191,115,171}\fi
  \fi}}
\null
\end{document}
'
%   lualatex -jobname=capa        '\def\qual{capa}% gerar-fundos.tex — gera fundo-impar.pdf, fundo-par.pdf e capa.pdf (A5) dos
% MWEs web e print. Arte provisória: faixas de cor, sem texto.
% Uso, uma vez por arquivo (o nome do PDF é o do -jobname):
%   lualatex -jobname=fundo-impar '\def\qual{impar}\input{gerar-fundos}'
%   lualatex -jobname=fundo-par   '\def\qual{par}\input{gerar-fundos}'
%   lualatex -jobname=capa        '\def\qual{capa}\input{gerar-fundos}'
% Os PDFs gerados ficam versionados ao lado (pequenos, vetoriais).
\documentclass{article}
\usepackage[paperwidth=148mm,paperheight=210mm,margin=0pt]{geometry}
\usepackage{xcolor,eso-pic}
\pagestyle{empty}
\def\faixa#1#2#3#4#5{\put(#1,#2){\color[RGB]{#5}\rule{#3}{#4}}}
\begin{document}
\AddToShipoutPictureBG*{\AtPageLowerLeft{%
  \expandafter\ifx\csname qual\endcsname\relax\else
  \def\a{impar}\ifx\qual\a   % faixa na margem externa à direita
    \faixa{130mm}{0mm}{18mm}{210mm}{217,238,242}\faixa{0mm}{0mm}{148mm}{3mm}{4,88,110}\fi
  \def\a{par}\ifx\qual\a     % faixa na margem externa à esquerda
    \faixa{0mm}{0mm}{18mm}{210mm}{217,238,242}\faixa{0mm}{0mm}{148mm}{3mm}{4,88,110}\fi
  \def\a{capa}\ifx\qual\a
    \faixa{0mm}{0mm}{148mm}{210mm}{25,159,175}\faixa{0mm}{150mm}{148mm}{20mm}{191,115,171}\fi
  \fi}}
\null
\end{document}
'
% Os PDFs gerados ficam versionados ao lado (pequenos, vetoriais).
\documentclass{article}
\usepackage[paperwidth=148mm,paperheight=210mm,margin=0pt]{geometry}
\usepackage{xcolor,eso-pic}
\pagestyle{empty}
\def\faixa#1#2#3#4#5{\put(#1,#2){\color[RGB]{#5}\rule{#3}{#4}}}
\begin{document}
\AddToShipoutPictureBG*{\AtPageLowerLeft{%
  \expandafter\ifx\csname qual\endcsname\relax\else
  \def\a{impar}\ifx\qual\a   % faixa na margem externa à direita
    \faixa{130mm}{0mm}{18mm}{210mm}{217,238,242}\faixa{0mm}{0mm}{148mm}{3mm}{4,88,110}\fi
  \def\a{par}\ifx\qual\a     % faixa na margem externa à esquerda
    \faixa{0mm}{0mm}{18mm}{210mm}{217,238,242}\faixa{0mm}{0mm}{148mm}{3mm}{4,88,110}\fi
  \def\a{capa}\ifx\qual\a
    \faixa{0mm}{0mm}{148mm}{210mm}{25,159,175}\faixa{0mm}{150mm}{148mm}{20mm}{191,115,171}\fi
  \fi}}
\null
\end{document}
'
%   lualatex -jobname=fundo-par   '\def\qual{par}% gerar-fundos.tex — gera fundo-impar.pdf, fundo-par.pdf e capa.pdf (A5) dos
% MWEs web e print. Arte provisória: faixas de cor, sem texto.
% Uso, uma vez por arquivo (o nome do PDF é o do -jobname):
%   lualatex -jobname=fundo-impar '\def\qual{impar}% gerar-fundos.tex — gera fundo-impar.pdf, fundo-par.pdf e capa.pdf (A5) dos
% MWEs web e print. Arte provisória: faixas de cor, sem texto.
% Uso, uma vez por arquivo (o nome do PDF é o do -jobname):
%   lualatex -jobname=fundo-impar '\def\qual{impar}\input{gerar-fundos}'
%   lualatex -jobname=fundo-par   '\def\qual{par}\input{gerar-fundos}'
%   lualatex -jobname=capa        '\def\qual{capa}\input{gerar-fundos}'
% Os PDFs gerados ficam versionados ao lado (pequenos, vetoriais).
\documentclass{article}
\usepackage[paperwidth=148mm,paperheight=210mm,margin=0pt]{geometry}
\usepackage{xcolor,eso-pic}
\pagestyle{empty}
\def\faixa#1#2#3#4#5{\put(#1,#2){\color[RGB]{#5}\rule{#3}{#4}}}
\begin{document}
\AddToShipoutPictureBG*{\AtPageLowerLeft{%
  \expandafter\ifx\csname qual\endcsname\relax\else
  \def\a{impar}\ifx\qual\a   % faixa na margem externa à direita
    \faixa{130mm}{0mm}{18mm}{210mm}{217,238,242}\faixa{0mm}{0mm}{148mm}{3mm}{4,88,110}\fi
  \def\a{par}\ifx\qual\a     % faixa na margem externa à esquerda
    \faixa{0mm}{0mm}{18mm}{210mm}{217,238,242}\faixa{0mm}{0mm}{148mm}{3mm}{4,88,110}\fi
  \def\a{capa}\ifx\qual\a
    \faixa{0mm}{0mm}{148mm}{210mm}{25,159,175}\faixa{0mm}{150mm}{148mm}{20mm}{191,115,171}\fi
  \fi}}
\null
\end{document}
'
%   lualatex -jobname=fundo-par   '\def\qual{par}% gerar-fundos.tex — gera fundo-impar.pdf, fundo-par.pdf e capa.pdf (A5) dos
% MWEs web e print. Arte provisória: faixas de cor, sem texto.
% Uso, uma vez por arquivo (o nome do PDF é o do -jobname):
%   lualatex -jobname=fundo-impar '\def\qual{impar}\input{gerar-fundos}'
%   lualatex -jobname=fundo-par   '\def\qual{par}\input{gerar-fundos}'
%   lualatex -jobname=capa        '\def\qual{capa}\input{gerar-fundos}'
% Os PDFs gerados ficam versionados ao lado (pequenos, vetoriais).
\documentclass{article}
\usepackage[paperwidth=148mm,paperheight=210mm,margin=0pt]{geometry}
\usepackage{xcolor,eso-pic}
\pagestyle{empty}
\def\faixa#1#2#3#4#5{\put(#1,#2){\color[RGB]{#5}\rule{#3}{#4}}}
\begin{document}
\AddToShipoutPictureBG*{\AtPageLowerLeft{%
  \expandafter\ifx\csname qual\endcsname\relax\else
  \def\a{impar}\ifx\qual\a   % faixa na margem externa à direita
    \faixa{130mm}{0mm}{18mm}{210mm}{217,238,242}\faixa{0mm}{0mm}{148mm}{3mm}{4,88,110}\fi
  \def\a{par}\ifx\qual\a     % faixa na margem externa à esquerda
    \faixa{0mm}{0mm}{18mm}{210mm}{217,238,242}\faixa{0mm}{0mm}{148mm}{3mm}{4,88,110}\fi
  \def\a{capa}\ifx\qual\a
    \faixa{0mm}{0mm}{148mm}{210mm}{25,159,175}\faixa{0mm}{150mm}{148mm}{20mm}{191,115,171}\fi
  \fi}}
\null
\end{document}
'
%   lualatex -jobname=capa        '\def\qual{capa}% gerar-fundos.tex — gera fundo-impar.pdf, fundo-par.pdf e capa.pdf (A5) dos
% MWEs web e print. Arte provisória: faixas de cor, sem texto.
% Uso, uma vez por arquivo (o nome do PDF é o do -jobname):
%   lualatex -jobname=fundo-impar '\def\qual{impar}\input{gerar-fundos}'
%   lualatex -jobname=fundo-par   '\def\qual{par}\input{gerar-fundos}'
%   lualatex -jobname=capa        '\def\qual{capa}\input{gerar-fundos}'
% Os PDFs gerados ficam versionados ao lado (pequenos, vetoriais).
\documentclass{article}
\usepackage[paperwidth=148mm,paperheight=210mm,margin=0pt]{geometry}
\usepackage{xcolor,eso-pic}
\pagestyle{empty}
\def\faixa#1#2#3#4#5{\put(#1,#2){\color[RGB]{#5}\rule{#3}{#4}}}
\begin{document}
\AddToShipoutPictureBG*{\AtPageLowerLeft{%
  \expandafter\ifx\csname qual\endcsname\relax\else
  \def\a{impar}\ifx\qual\a   % faixa na margem externa à direita
    \faixa{130mm}{0mm}{18mm}{210mm}{217,238,242}\faixa{0mm}{0mm}{148mm}{3mm}{4,88,110}\fi
  \def\a{par}\ifx\qual\a     % faixa na margem externa à esquerda
    \faixa{0mm}{0mm}{18mm}{210mm}{217,238,242}\faixa{0mm}{0mm}{148mm}{3mm}{4,88,110}\fi
  \def\a{capa}\ifx\qual\a
    \faixa{0mm}{0mm}{148mm}{210mm}{25,159,175}\faixa{0mm}{150mm}{148mm}{20mm}{191,115,171}\fi
  \fi}}
\null
\end{document}
'
% Os PDFs gerados ficam versionados ao lado (pequenos, vetoriais).
\documentclass{article}
\usepackage[paperwidth=148mm,paperheight=210mm,margin=0pt]{geometry}
\usepackage{xcolor,eso-pic}
\pagestyle{empty}
\def\faixa#1#2#3#4#5{\put(#1,#2){\color[RGB]{#5}\rule{#3}{#4}}}
\begin{document}
\AddToShipoutPictureBG*{\AtPageLowerLeft{%
  \expandafter\ifx\csname qual\endcsname\relax\else
  \def\a{impar}\ifx\qual\a   % faixa na margem externa à direita
    \faixa{130mm}{0mm}{18mm}{210mm}{217,238,242}\faixa{0mm}{0mm}{148mm}{3mm}{4,88,110}\fi
  \def\a{par}\ifx\qual\a     % faixa na margem externa à esquerda
    \faixa{0mm}{0mm}{18mm}{210mm}{217,238,242}\faixa{0mm}{0mm}{148mm}{3mm}{4,88,110}\fi
  \def\a{capa}\ifx\qual\a
    \faixa{0mm}{0mm}{148mm}{210mm}{25,159,175}\faixa{0mm}{150mm}{148mm}{20mm}{191,115,171}\fi
  \fi}}
\null
\end{document}
'
%   lualatex -jobname=capa        '\def\qual{capa}% gerar-fundos.tex — gera fundo-impar.pdf, fundo-par.pdf e capa.pdf (A5) dos
% MWEs web e print. Arte provisória: faixas de cor, sem texto.
% Uso, uma vez por arquivo (o nome do PDF é o do -jobname):
%   lualatex -jobname=fundo-impar '\def\qual{impar}% gerar-fundos.tex — gera fundo-impar.pdf, fundo-par.pdf e capa.pdf (A5) dos
% MWEs web e print. Arte provisória: faixas de cor, sem texto.
% Uso, uma vez por arquivo (o nome do PDF é o do -jobname):
%   lualatex -jobname=fundo-impar '\def\qual{impar}\input{gerar-fundos}'
%   lualatex -jobname=fundo-par   '\def\qual{par}\input{gerar-fundos}'
%   lualatex -jobname=capa        '\def\qual{capa}\input{gerar-fundos}'
% Os PDFs gerados ficam versionados ao lado (pequenos, vetoriais).
\documentclass{article}
\usepackage[paperwidth=148mm,paperheight=210mm,margin=0pt]{geometry}
\usepackage{xcolor,eso-pic}
\pagestyle{empty}
\def\faixa#1#2#3#4#5{\put(#1,#2){\color[RGB]{#5}\rule{#3}{#4}}}
\begin{document}
\AddToShipoutPictureBG*{\AtPageLowerLeft{%
  \expandafter\ifx\csname qual\endcsname\relax\else
  \def\a{impar}\ifx\qual\a   % faixa na margem externa à direita
    \faixa{130mm}{0mm}{18mm}{210mm}{217,238,242}\faixa{0mm}{0mm}{148mm}{3mm}{4,88,110}\fi
  \def\a{par}\ifx\qual\a     % faixa na margem externa à esquerda
    \faixa{0mm}{0mm}{18mm}{210mm}{217,238,242}\faixa{0mm}{0mm}{148mm}{3mm}{4,88,110}\fi
  \def\a{capa}\ifx\qual\a
    \faixa{0mm}{0mm}{148mm}{210mm}{25,159,175}\faixa{0mm}{150mm}{148mm}{20mm}{191,115,171}\fi
  \fi}}
\null
\end{document}
'
%   lualatex -jobname=fundo-par   '\def\qual{par}% gerar-fundos.tex — gera fundo-impar.pdf, fundo-par.pdf e capa.pdf (A5) dos
% MWEs web e print. Arte provisória: faixas de cor, sem texto.
% Uso, uma vez por arquivo (o nome do PDF é o do -jobname):
%   lualatex -jobname=fundo-impar '\def\qual{impar}\input{gerar-fundos}'
%   lualatex -jobname=fundo-par   '\def\qual{par}\input{gerar-fundos}'
%   lualatex -jobname=capa        '\def\qual{capa}\input{gerar-fundos}'
% Os PDFs gerados ficam versionados ao lado (pequenos, vetoriais).
\documentclass{article}
\usepackage[paperwidth=148mm,paperheight=210mm,margin=0pt]{geometry}
\usepackage{xcolor,eso-pic}
\pagestyle{empty}
\def\faixa#1#2#3#4#5{\put(#1,#2){\color[RGB]{#5}\rule{#3}{#4}}}
\begin{document}
\AddToShipoutPictureBG*{\AtPageLowerLeft{%
  \expandafter\ifx\csname qual\endcsname\relax\else
  \def\a{impar}\ifx\qual\a   % faixa na margem externa à direita
    \faixa{130mm}{0mm}{18mm}{210mm}{217,238,242}\faixa{0mm}{0mm}{148mm}{3mm}{4,88,110}\fi
  \def\a{par}\ifx\qual\a     % faixa na margem externa à esquerda
    \faixa{0mm}{0mm}{18mm}{210mm}{217,238,242}\faixa{0mm}{0mm}{148mm}{3mm}{4,88,110}\fi
  \def\a{capa}\ifx\qual\a
    \faixa{0mm}{0mm}{148mm}{210mm}{25,159,175}\faixa{0mm}{150mm}{148mm}{20mm}{191,115,171}\fi
  \fi}}
\null
\end{document}
'
%   lualatex -jobname=capa        '\def\qual{capa}% gerar-fundos.tex — gera fundo-impar.pdf, fundo-par.pdf e capa.pdf (A5) dos
% MWEs web e print. Arte provisória: faixas de cor, sem texto.
% Uso, uma vez por arquivo (o nome do PDF é o do -jobname):
%   lualatex -jobname=fundo-impar '\def\qual{impar}\input{gerar-fundos}'
%   lualatex -jobname=fundo-par   '\def\qual{par}\input{gerar-fundos}'
%   lualatex -jobname=capa        '\def\qual{capa}\input{gerar-fundos}'
% Os PDFs gerados ficam versionados ao lado (pequenos, vetoriais).
\documentclass{article}
\usepackage[paperwidth=148mm,paperheight=210mm,margin=0pt]{geometry}
\usepackage{xcolor,eso-pic}
\pagestyle{empty}
\def\faixa#1#2#3#4#5{\put(#1,#2){\color[RGB]{#5}\rule{#3}{#4}}}
\begin{document}
\AddToShipoutPictureBG*{\AtPageLowerLeft{%
  \expandafter\ifx\csname qual\endcsname\relax\else
  \def\a{impar}\ifx\qual\a   % faixa na margem externa à direita
    \faixa{130mm}{0mm}{18mm}{210mm}{217,238,242}\faixa{0mm}{0mm}{148mm}{3mm}{4,88,110}\fi
  \def\a{par}\ifx\qual\a     % faixa na margem externa à esquerda
    \faixa{0mm}{0mm}{18mm}{210mm}{217,238,242}\faixa{0mm}{0mm}{148mm}{3mm}{4,88,110}\fi
  \def\a{capa}\ifx\qual\a
    \faixa{0mm}{0mm}{148mm}{210mm}{25,159,175}\faixa{0mm}{150mm}{148mm}{20mm}{191,115,171}\fi
  \fi}}
\null
\end{document}
'
% Os PDFs gerados ficam versionados ao lado (pequenos, vetoriais).
\documentclass{article}
\usepackage[paperwidth=148mm,paperheight=210mm,margin=0pt]{geometry}
\usepackage{xcolor,eso-pic}
\pagestyle{empty}
\def\faixa#1#2#3#4#5{\put(#1,#2){\color[RGB]{#5}\rule{#3}{#4}}}
\begin{document}
\AddToShipoutPictureBG*{\AtPageLowerLeft{%
  \expandafter\ifx\csname qual\endcsname\relax\else
  \def\a{impar}\ifx\qual\a   % faixa na margem externa à direita
    \faixa{130mm}{0mm}{18mm}{210mm}{217,238,242}\faixa{0mm}{0mm}{148mm}{3mm}{4,88,110}\fi
  \def\a{par}\ifx\qual\a     % faixa na margem externa à esquerda
    \faixa{0mm}{0mm}{18mm}{210mm}{217,238,242}\faixa{0mm}{0mm}{148mm}{3mm}{4,88,110}\fi
  \def\a{capa}\ifx\qual\a
    \faixa{0mm}{0mm}{148mm}{210mm}{25,159,175}\faixa{0mm}{150mm}{148mm}{20mm}{191,115,171}\fi
  \fi}}
\null
\end{document}
'
% Os PDFs gerados ficam versionados ao lado (pequenos, vetoriais).
\documentclass{article}
\usepackage[paperwidth=148mm,paperheight=210mm,margin=0pt]{geometry}
\usepackage{xcolor,eso-pic}
\pagestyle{empty}
\def\faixa#1#2#3#4#5{\put(#1,#2){\color[RGB]{#5}\rule{#3}{#4}}}
\begin{document}
\AddToShipoutPictureBG*{\AtPageLowerLeft{%
  \expandafter\ifx\csname qual\endcsname\relax\else
  \def\a{impar}\ifx\qual\a   % faixa na margem externa à direita
    \faixa{130mm}{0mm}{18mm}{210mm}{217,238,242}\faixa{0mm}{0mm}{148mm}{3mm}{4,88,110}\fi
  \def\a{par}\ifx\qual\a     % faixa na margem externa à esquerda
    \faixa{0mm}{0mm}{18mm}{210mm}{217,238,242}\faixa{0mm}{0mm}{148mm}{3mm}{4,88,110}\fi
  \def\a{capa}\ifx\qual\a
    \faixa{0mm}{0mm}{148mm}{210mm}{25,159,175}\faixa{0mm}{150mm}{148mm}{20mm}{191,115,171}\fi
  \fi}}
\null
\end{document}
'
%   lualatex -jobname=capa        '\def\qual{capa}% gerar-fundos.tex — gera fundo-impar.pdf, fundo-par.pdf e capa.pdf (A5) dos
% MWEs web e print. Arte provisória: faixas de cor, sem texto.
% Uso, uma vez por arquivo (o nome do PDF é o do -jobname):
%   lualatex -jobname=fundo-impar '\def\qual{impar}% gerar-fundos.tex — gera fundo-impar.pdf, fundo-par.pdf e capa.pdf (A5) dos
% MWEs web e print. Arte provisória: faixas de cor, sem texto.
% Uso, uma vez por arquivo (o nome do PDF é o do -jobname):
%   lualatex -jobname=fundo-impar '\def\qual{impar}% gerar-fundos.tex — gera fundo-impar.pdf, fundo-par.pdf e capa.pdf (A5) dos
% MWEs web e print. Arte provisória: faixas de cor, sem texto.
% Uso, uma vez por arquivo (o nome do PDF é o do -jobname):
%   lualatex -jobname=fundo-impar '\def\qual{impar}\input{gerar-fundos}'
%   lualatex -jobname=fundo-par   '\def\qual{par}\input{gerar-fundos}'
%   lualatex -jobname=capa        '\def\qual{capa}\input{gerar-fundos}'
% Os PDFs gerados ficam versionados ao lado (pequenos, vetoriais).
\documentclass{article}
\usepackage[paperwidth=148mm,paperheight=210mm,margin=0pt]{geometry}
\usepackage{xcolor,eso-pic}
\pagestyle{empty}
\def\faixa#1#2#3#4#5{\put(#1,#2){\color[RGB]{#5}\rule{#3}{#4}}}
\begin{document}
\AddToShipoutPictureBG*{\AtPageLowerLeft{%
  \expandafter\ifx\csname qual\endcsname\relax\else
  \def\a{impar}\ifx\qual\a   % faixa na margem externa à direita
    \faixa{130mm}{0mm}{18mm}{210mm}{217,238,242}\faixa{0mm}{0mm}{148mm}{3mm}{4,88,110}\fi
  \def\a{par}\ifx\qual\a     % faixa na margem externa à esquerda
    \faixa{0mm}{0mm}{18mm}{210mm}{217,238,242}\faixa{0mm}{0mm}{148mm}{3mm}{4,88,110}\fi
  \def\a{capa}\ifx\qual\a
    \faixa{0mm}{0mm}{148mm}{210mm}{25,159,175}\faixa{0mm}{150mm}{148mm}{20mm}{191,115,171}\fi
  \fi}}
\null
\end{document}
'
%   lualatex -jobname=fundo-par   '\def\qual{par}% gerar-fundos.tex — gera fundo-impar.pdf, fundo-par.pdf e capa.pdf (A5) dos
% MWEs web e print. Arte provisória: faixas de cor, sem texto.
% Uso, uma vez por arquivo (o nome do PDF é o do -jobname):
%   lualatex -jobname=fundo-impar '\def\qual{impar}\input{gerar-fundos}'
%   lualatex -jobname=fundo-par   '\def\qual{par}\input{gerar-fundos}'
%   lualatex -jobname=capa        '\def\qual{capa}\input{gerar-fundos}'
% Os PDFs gerados ficam versionados ao lado (pequenos, vetoriais).
\documentclass{article}
\usepackage[paperwidth=148mm,paperheight=210mm,margin=0pt]{geometry}
\usepackage{xcolor,eso-pic}
\pagestyle{empty}
\def\faixa#1#2#3#4#5{\put(#1,#2){\color[RGB]{#5}\rule{#3}{#4}}}
\begin{document}
\AddToShipoutPictureBG*{\AtPageLowerLeft{%
  \expandafter\ifx\csname qual\endcsname\relax\else
  \def\a{impar}\ifx\qual\a   % faixa na margem externa à direita
    \faixa{130mm}{0mm}{18mm}{210mm}{217,238,242}\faixa{0mm}{0mm}{148mm}{3mm}{4,88,110}\fi
  \def\a{par}\ifx\qual\a     % faixa na margem externa à esquerda
    \faixa{0mm}{0mm}{18mm}{210mm}{217,238,242}\faixa{0mm}{0mm}{148mm}{3mm}{4,88,110}\fi
  \def\a{capa}\ifx\qual\a
    \faixa{0mm}{0mm}{148mm}{210mm}{25,159,175}\faixa{0mm}{150mm}{148mm}{20mm}{191,115,171}\fi
  \fi}}
\null
\end{document}
'
%   lualatex -jobname=capa        '\def\qual{capa}% gerar-fundos.tex — gera fundo-impar.pdf, fundo-par.pdf e capa.pdf (A5) dos
% MWEs web e print. Arte provisória: faixas de cor, sem texto.
% Uso, uma vez por arquivo (o nome do PDF é o do -jobname):
%   lualatex -jobname=fundo-impar '\def\qual{impar}\input{gerar-fundos}'
%   lualatex -jobname=fundo-par   '\def\qual{par}\input{gerar-fundos}'
%   lualatex -jobname=capa        '\def\qual{capa}\input{gerar-fundos}'
% Os PDFs gerados ficam versionados ao lado (pequenos, vetoriais).
\documentclass{article}
\usepackage[paperwidth=148mm,paperheight=210mm,margin=0pt]{geometry}
\usepackage{xcolor,eso-pic}
\pagestyle{empty}
\def\faixa#1#2#3#4#5{\put(#1,#2){\color[RGB]{#5}\rule{#3}{#4}}}
\begin{document}
\AddToShipoutPictureBG*{\AtPageLowerLeft{%
  \expandafter\ifx\csname qual\endcsname\relax\else
  \def\a{impar}\ifx\qual\a   % faixa na margem externa à direita
    \faixa{130mm}{0mm}{18mm}{210mm}{217,238,242}\faixa{0mm}{0mm}{148mm}{3mm}{4,88,110}\fi
  \def\a{par}\ifx\qual\a     % faixa na margem externa à esquerda
    \faixa{0mm}{0mm}{18mm}{210mm}{217,238,242}\faixa{0mm}{0mm}{148mm}{3mm}{4,88,110}\fi
  \def\a{capa}\ifx\qual\a
    \faixa{0mm}{0mm}{148mm}{210mm}{25,159,175}\faixa{0mm}{150mm}{148mm}{20mm}{191,115,171}\fi
  \fi}}
\null
\end{document}
'
% Os PDFs gerados ficam versionados ao lado (pequenos, vetoriais).
\documentclass{article}
\usepackage[paperwidth=148mm,paperheight=210mm,margin=0pt]{geometry}
\usepackage{xcolor,eso-pic}
\pagestyle{empty}
\def\faixa#1#2#3#4#5{\put(#1,#2){\color[RGB]{#5}\rule{#3}{#4}}}
\begin{document}
\AddToShipoutPictureBG*{\AtPageLowerLeft{%
  \expandafter\ifx\csname qual\endcsname\relax\else
  \def\a{impar}\ifx\qual\a   % faixa na margem externa à direita
    \faixa{130mm}{0mm}{18mm}{210mm}{217,238,242}\faixa{0mm}{0mm}{148mm}{3mm}{4,88,110}\fi
  \def\a{par}\ifx\qual\a     % faixa na margem externa à esquerda
    \faixa{0mm}{0mm}{18mm}{210mm}{217,238,242}\faixa{0mm}{0mm}{148mm}{3mm}{4,88,110}\fi
  \def\a{capa}\ifx\qual\a
    \faixa{0mm}{0mm}{148mm}{210mm}{25,159,175}\faixa{0mm}{150mm}{148mm}{20mm}{191,115,171}\fi
  \fi}}
\null
\end{document}
'
%   lualatex -jobname=fundo-par   '\def\qual{par}% gerar-fundos.tex — gera fundo-impar.pdf, fundo-par.pdf e capa.pdf (A5) dos
% MWEs web e print. Arte provisória: faixas de cor, sem texto.
% Uso, uma vez por arquivo (o nome do PDF é o do -jobname):
%   lualatex -jobname=fundo-impar '\def\qual{impar}% gerar-fundos.tex — gera fundo-impar.pdf, fundo-par.pdf e capa.pdf (A5) dos
% MWEs web e print. Arte provisória: faixas de cor, sem texto.
% Uso, uma vez por arquivo (o nome do PDF é o do -jobname):
%   lualatex -jobname=fundo-impar '\def\qual{impar}\input{gerar-fundos}'
%   lualatex -jobname=fundo-par   '\def\qual{par}\input{gerar-fundos}'
%   lualatex -jobname=capa        '\def\qual{capa}\input{gerar-fundos}'
% Os PDFs gerados ficam versionados ao lado (pequenos, vetoriais).
\documentclass{article}
\usepackage[paperwidth=148mm,paperheight=210mm,margin=0pt]{geometry}
\usepackage{xcolor,eso-pic}
\pagestyle{empty}
\def\faixa#1#2#3#4#5{\put(#1,#2){\color[RGB]{#5}\rule{#3}{#4}}}
\begin{document}
\AddToShipoutPictureBG*{\AtPageLowerLeft{%
  \expandafter\ifx\csname qual\endcsname\relax\else
  \def\a{impar}\ifx\qual\a   % faixa na margem externa à direita
    \faixa{130mm}{0mm}{18mm}{210mm}{217,238,242}\faixa{0mm}{0mm}{148mm}{3mm}{4,88,110}\fi
  \def\a{par}\ifx\qual\a     % faixa na margem externa à esquerda
    \faixa{0mm}{0mm}{18mm}{210mm}{217,238,242}\faixa{0mm}{0mm}{148mm}{3mm}{4,88,110}\fi
  \def\a{capa}\ifx\qual\a
    \faixa{0mm}{0mm}{148mm}{210mm}{25,159,175}\faixa{0mm}{150mm}{148mm}{20mm}{191,115,171}\fi
  \fi}}
\null
\end{document}
'
%   lualatex -jobname=fundo-par   '\def\qual{par}% gerar-fundos.tex — gera fundo-impar.pdf, fundo-par.pdf e capa.pdf (A5) dos
% MWEs web e print. Arte provisória: faixas de cor, sem texto.
% Uso, uma vez por arquivo (o nome do PDF é o do -jobname):
%   lualatex -jobname=fundo-impar '\def\qual{impar}\input{gerar-fundos}'
%   lualatex -jobname=fundo-par   '\def\qual{par}\input{gerar-fundos}'
%   lualatex -jobname=capa        '\def\qual{capa}\input{gerar-fundos}'
% Os PDFs gerados ficam versionados ao lado (pequenos, vetoriais).
\documentclass{article}
\usepackage[paperwidth=148mm,paperheight=210mm,margin=0pt]{geometry}
\usepackage{xcolor,eso-pic}
\pagestyle{empty}
\def\faixa#1#2#3#4#5{\put(#1,#2){\color[RGB]{#5}\rule{#3}{#4}}}
\begin{document}
\AddToShipoutPictureBG*{\AtPageLowerLeft{%
  \expandafter\ifx\csname qual\endcsname\relax\else
  \def\a{impar}\ifx\qual\a   % faixa na margem externa à direita
    \faixa{130mm}{0mm}{18mm}{210mm}{217,238,242}\faixa{0mm}{0mm}{148mm}{3mm}{4,88,110}\fi
  \def\a{par}\ifx\qual\a     % faixa na margem externa à esquerda
    \faixa{0mm}{0mm}{18mm}{210mm}{217,238,242}\faixa{0mm}{0mm}{148mm}{3mm}{4,88,110}\fi
  \def\a{capa}\ifx\qual\a
    \faixa{0mm}{0mm}{148mm}{210mm}{25,159,175}\faixa{0mm}{150mm}{148mm}{20mm}{191,115,171}\fi
  \fi}}
\null
\end{document}
'
%   lualatex -jobname=capa        '\def\qual{capa}% gerar-fundos.tex — gera fundo-impar.pdf, fundo-par.pdf e capa.pdf (A5) dos
% MWEs web e print. Arte provisória: faixas de cor, sem texto.
% Uso, uma vez por arquivo (o nome do PDF é o do -jobname):
%   lualatex -jobname=fundo-impar '\def\qual{impar}\input{gerar-fundos}'
%   lualatex -jobname=fundo-par   '\def\qual{par}\input{gerar-fundos}'
%   lualatex -jobname=capa        '\def\qual{capa}\input{gerar-fundos}'
% Os PDFs gerados ficam versionados ao lado (pequenos, vetoriais).
\documentclass{article}
\usepackage[paperwidth=148mm,paperheight=210mm,margin=0pt]{geometry}
\usepackage{xcolor,eso-pic}
\pagestyle{empty}
\def\faixa#1#2#3#4#5{\put(#1,#2){\color[RGB]{#5}\rule{#3}{#4}}}
\begin{document}
\AddToShipoutPictureBG*{\AtPageLowerLeft{%
  \expandafter\ifx\csname qual\endcsname\relax\else
  \def\a{impar}\ifx\qual\a   % faixa na margem externa à direita
    \faixa{130mm}{0mm}{18mm}{210mm}{217,238,242}\faixa{0mm}{0mm}{148mm}{3mm}{4,88,110}\fi
  \def\a{par}\ifx\qual\a     % faixa na margem externa à esquerda
    \faixa{0mm}{0mm}{18mm}{210mm}{217,238,242}\faixa{0mm}{0mm}{148mm}{3mm}{4,88,110}\fi
  \def\a{capa}\ifx\qual\a
    \faixa{0mm}{0mm}{148mm}{210mm}{25,159,175}\faixa{0mm}{150mm}{148mm}{20mm}{191,115,171}\fi
  \fi}}
\null
\end{document}
'
% Os PDFs gerados ficam versionados ao lado (pequenos, vetoriais).
\documentclass{article}
\usepackage[paperwidth=148mm,paperheight=210mm,margin=0pt]{geometry}
\usepackage{xcolor,eso-pic}
\pagestyle{empty}
\def\faixa#1#2#3#4#5{\put(#1,#2){\color[RGB]{#5}\rule{#3}{#4}}}
\begin{document}
\AddToShipoutPictureBG*{\AtPageLowerLeft{%
  \expandafter\ifx\csname qual\endcsname\relax\else
  \def\a{impar}\ifx\qual\a   % faixa na margem externa à direita
    \faixa{130mm}{0mm}{18mm}{210mm}{217,238,242}\faixa{0mm}{0mm}{148mm}{3mm}{4,88,110}\fi
  \def\a{par}\ifx\qual\a     % faixa na margem externa à esquerda
    \faixa{0mm}{0mm}{18mm}{210mm}{217,238,242}\faixa{0mm}{0mm}{148mm}{3mm}{4,88,110}\fi
  \def\a{capa}\ifx\qual\a
    \faixa{0mm}{0mm}{148mm}{210mm}{25,159,175}\faixa{0mm}{150mm}{148mm}{20mm}{191,115,171}\fi
  \fi}}
\null
\end{document}
'
%   lualatex -jobname=capa        '\def\qual{capa}% gerar-fundos.tex — gera fundo-impar.pdf, fundo-par.pdf e capa.pdf (A5) dos
% MWEs web e print. Arte provisória: faixas de cor, sem texto.
% Uso, uma vez por arquivo (o nome do PDF é o do -jobname):
%   lualatex -jobname=fundo-impar '\def\qual{impar}% gerar-fundos.tex — gera fundo-impar.pdf, fundo-par.pdf e capa.pdf (A5) dos
% MWEs web e print. Arte provisória: faixas de cor, sem texto.
% Uso, uma vez por arquivo (o nome do PDF é o do -jobname):
%   lualatex -jobname=fundo-impar '\def\qual{impar}\input{gerar-fundos}'
%   lualatex -jobname=fundo-par   '\def\qual{par}\input{gerar-fundos}'
%   lualatex -jobname=capa        '\def\qual{capa}\input{gerar-fundos}'
% Os PDFs gerados ficam versionados ao lado (pequenos, vetoriais).
\documentclass{article}
\usepackage[paperwidth=148mm,paperheight=210mm,margin=0pt]{geometry}
\usepackage{xcolor,eso-pic}
\pagestyle{empty}
\def\faixa#1#2#3#4#5{\put(#1,#2){\color[RGB]{#5}\rule{#3}{#4}}}
\begin{document}
\AddToShipoutPictureBG*{\AtPageLowerLeft{%
  \expandafter\ifx\csname qual\endcsname\relax\else
  \def\a{impar}\ifx\qual\a   % faixa na margem externa à direita
    \faixa{130mm}{0mm}{18mm}{210mm}{217,238,242}\faixa{0mm}{0mm}{148mm}{3mm}{4,88,110}\fi
  \def\a{par}\ifx\qual\a     % faixa na margem externa à esquerda
    \faixa{0mm}{0mm}{18mm}{210mm}{217,238,242}\faixa{0mm}{0mm}{148mm}{3mm}{4,88,110}\fi
  \def\a{capa}\ifx\qual\a
    \faixa{0mm}{0mm}{148mm}{210mm}{25,159,175}\faixa{0mm}{150mm}{148mm}{20mm}{191,115,171}\fi
  \fi}}
\null
\end{document}
'
%   lualatex -jobname=fundo-par   '\def\qual{par}% gerar-fundos.tex — gera fundo-impar.pdf, fundo-par.pdf e capa.pdf (A5) dos
% MWEs web e print. Arte provisória: faixas de cor, sem texto.
% Uso, uma vez por arquivo (o nome do PDF é o do -jobname):
%   lualatex -jobname=fundo-impar '\def\qual{impar}\input{gerar-fundos}'
%   lualatex -jobname=fundo-par   '\def\qual{par}\input{gerar-fundos}'
%   lualatex -jobname=capa        '\def\qual{capa}\input{gerar-fundos}'
% Os PDFs gerados ficam versionados ao lado (pequenos, vetoriais).
\documentclass{article}
\usepackage[paperwidth=148mm,paperheight=210mm,margin=0pt]{geometry}
\usepackage{xcolor,eso-pic}
\pagestyle{empty}
\def\faixa#1#2#3#4#5{\put(#1,#2){\color[RGB]{#5}\rule{#3}{#4}}}
\begin{document}
\AddToShipoutPictureBG*{\AtPageLowerLeft{%
  \expandafter\ifx\csname qual\endcsname\relax\else
  \def\a{impar}\ifx\qual\a   % faixa na margem externa à direita
    \faixa{130mm}{0mm}{18mm}{210mm}{217,238,242}\faixa{0mm}{0mm}{148mm}{3mm}{4,88,110}\fi
  \def\a{par}\ifx\qual\a     % faixa na margem externa à esquerda
    \faixa{0mm}{0mm}{18mm}{210mm}{217,238,242}\faixa{0mm}{0mm}{148mm}{3mm}{4,88,110}\fi
  \def\a{capa}\ifx\qual\a
    \faixa{0mm}{0mm}{148mm}{210mm}{25,159,175}\faixa{0mm}{150mm}{148mm}{20mm}{191,115,171}\fi
  \fi}}
\null
\end{document}
'
%   lualatex -jobname=capa        '\def\qual{capa}% gerar-fundos.tex — gera fundo-impar.pdf, fundo-par.pdf e capa.pdf (A5) dos
% MWEs web e print. Arte provisória: faixas de cor, sem texto.
% Uso, uma vez por arquivo (o nome do PDF é o do -jobname):
%   lualatex -jobname=fundo-impar '\def\qual{impar}\input{gerar-fundos}'
%   lualatex -jobname=fundo-par   '\def\qual{par}\input{gerar-fundos}'
%   lualatex -jobname=capa        '\def\qual{capa}\input{gerar-fundos}'
% Os PDFs gerados ficam versionados ao lado (pequenos, vetoriais).
\documentclass{article}
\usepackage[paperwidth=148mm,paperheight=210mm,margin=0pt]{geometry}
\usepackage{xcolor,eso-pic}
\pagestyle{empty}
\def\faixa#1#2#3#4#5{\put(#1,#2){\color[RGB]{#5}\rule{#3}{#4}}}
\begin{document}
\AddToShipoutPictureBG*{\AtPageLowerLeft{%
  \expandafter\ifx\csname qual\endcsname\relax\else
  \def\a{impar}\ifx\qual\a   % faixa na margem externa à direita
    \faixa{130mm}{0mm}{18mm}{210mm}{217,238,242}\faixa{0mm}{0mm}{148mm}{3mm}{4,88,110}\fi
  \def\a{par}\ifx\qual\a     % faixa na margem externa à esquerda
    \faixa{0mm}{0mm}{18mm}{210mm}{217,238,242}\faixa{0mm}{0mm}{148mm}{3mm}{4,88,110}\fi
  \def\a{capa}\ifx\qual\a
    \faixa{0mm}{0mm}{148mm}{210mm}{25,159,175}\faixa{0mm}{150mm}{148mm}{20mm}{191,115,171}\fi
  \fi}}
\null
\end{document}
'
% Os PDFs gerados ficam versionados ao lado (pequenos, vetoriais).
\documentclass{article}
\usepackage[paperwidth=148mm,paperheight=210mm,margin=0pt]{geometry}
\usepackage{xcolor,eso-pic}
\pagestyle{empty}
\def\faixa#1#2#3#4#5{\put(#1,#2){\color[RGB]{#5}\rule{#3}{#4}}}
\begin{document}
\AddToShipoutPictureBG*{\AtPageLowerLeft{%
  \expandafter\ifx\csname qual\endcsname\relax\else
  \def\a{impar}\ifx\qual\a   % faixa na margem externa à direita
    \faixa{130mm}{0mm}{18mm}{210mm}{217,238,242}\faixa{0mm}{0mm}{148mm}{3mm}{4,88,110}\fi
  \def\a{par}\ifx\qual\a     % faixa na margem externa à esquerda
    \faixa{0mm}{0mm}{18mm}{210mm}{217,238,242}\faixa{0mm}{0mm}{148mm}{3mm}{4,88,110}\fi
  \def\a{capa}\ifx\qual\a
    \faixa{0mm}{0mm}{148mm}{210mm}{25,159,175}\faixa{0mm}{150mm}{148mm}{20mm}{191,115,171}\fi
  \fi}}
\null
\end{document}
'
% Os PDFs gerados ficam versionados ao lado (pequenos, vetoriais).
\documentclass{article}
\usepackage[paperwidth=148mm,paperheight=210mm,margin=0pt]{geometry}
\usepackage{xcolor,eso-pic}
\pagestyle{empty}
\def\faixa#1#2#3#4#5{\put(#1,#2){\color[RGB]{#5}\rule{#3}{#4}}}
\begin{document}
\AddToShipoutPictureBG*{\AtPageLowerLeft{%
  \expandafter\ifx\csname qual\endcsname\relax\else
  \def\a{impar}\ifx\qual\a   % faixa na margem externa à direita
    \faixa{130mm}{0mm}{18mm}{210mm}{217,238,242}\faixa{0mm}{0mm}{148mm}{3mm}{4,88,110}\fi
  \def\a{par}\ifx\qual\a     % faixa na margem externa à esquerda
    \faixa{0mm}{0mm}{18mm}{210mm}{217,238,242}\faixa{0mm}{0mm}{148mm}{3mm}{4,88,110}\fi
  \def\a{capa}\ifx\qual\a
    \faixa{0mm}{0mm}{148mm}{210mm}{25,159,175}\faixa{0mm}{150mm}{148mm}{20mm}{191,115,171}\fi
  \fi}}
\null
\end{document}
'
% Os PDFs gerados ficam versionados ao lado (pequenos, vetoriais).
\documentclass{article}
\usepackage[paperwidth=148mm,paperheight=210mm,margin=0pt]{geometry}
\usepackage{xcolor,eso-pic}
\pagestyle{empty}
\def\faixa#1#2#3#4#5{\put(#1,#2){\color[RGB]{#5}\rule{#3}{#4}}}
\begin{document}
\AddToShipoutPictureBG*{\AtPageLowerLeft{%
  \expandafter\ifx\csname qual\endcsname\relax\else
  \def\a{impar}\ifx\qual\a   % faixa na margem externa à direita
    \faixa{130mm}{0mm}{18mm}{210mm}{217,238,242}\faixa{0mm}{0mm}{148mm}{3mm}{4,88,110}\fi
  \def\a{par}\ifx\qual\a     % faixa na margem externa à esquerda
    \faixa{0mm}{0mm}{18mm}{210mm}{217,238,242}\faixa{0mm}{0mm}{148mm}{3mm}{4,88,110}\fi
  \def\a{capa}\ifx\qual\a
    \faixa{0mm}{0mm}{148mm}{210mm}{25,159,175}\faixa{0mm}{150mm}{148mm}{20mm}{191,115,171}\fi
  \fi}}
\null
\end{document}
